\documentclass[12pt]{article}
\usepackage[a4paper,margin=1in]{geometry}
\usepackage[utf8]{inputenc}
\usepackage[T1]{fontenc}
\usepackage{amsmath,amssymb}
\usepackage{textcomp}
\usepackage{parskip}
\begin{document}
\begin{center}
{\Large \textbf{2021 STEP III}}
\end{center}

\subsection*{Page 1}
\begin{flushleft}
\bigskip\textbf{Section A: Pure Mathematics}\par
\bigskip\textbf{Question 1.} (i) A curve has parametric equations\par
\[x = -4 cos 3 t , y = 12 sin t - 4 sin 3 t .\]
Find the equation of the normal to this curve at the point\par
(-4 cos 3 φ , 12 sin φ - 4 sin 3 φ),\par
\[where 0 < φ < 1 \pi.\]
Verify that this normal is a tangent to the curve\par
\[x 2 3 + y 2 3 = 4\]
at the point (8 cos 3 φ , 8 sin 3 φ).\par
\textbf{(ii)} A curve has parametric equations\par
\[x = cos t + t sin t , y = sin t - t cos t.\]
Find the equation of the normal to this curve at the point\par
(cos φ + φ sin φ , sin φ - φ cos φ),\par
\[where 0 < φ < 1 \pi.\]
Determine the perpendicular distance from the origin to this normal, and hence find\par
the equation of a curve, independent of φ, to which this normal is a tangent.\par
\end{flushleft}
\subsection*{Page 2}
\begin{flushleft}
\bigskip\textbf{Question 2.} (i) Let\par
a b c\par
\[x = , y = and z = ,\]
b - c c - a a - b\par
where a, b and c are distinct real numbers.\par
Show that ⎛ ⎞ ⎛ ⎞ ⎛ ⎞\par
1 -x x a 0\par
\[⎝ y 1 -y ⎠ ⎝ b ⎠ = ⎝ 0 ⎠\]
-z z 1 c 0\par
\[and use this result to deduce that yz + zx + xy = -1 .\]
Hence show that\par
\[a^2 b^2 c^2\]
\[+ + \leq 2 .\]
(b - c)2 (c - a)2 (a - b)2\par
\textbf{(ii)} Let\par
2a 2b 2c\par
\[x = , y = and z = ,\]
b + c c + a a + b\par
where a, b and c are positive real numbers.\par
\[Using a suitable matrix, show that xyz + yz + zx + xy = 4 .\]
Hence show that\par
\[(2a + b + c)(a + 2b + c)(a + b + 2c) > 5(b + c)(c + a)(a + b) .\]
Show further that\par
\[(2a + b + c)(a + 2b + c)(a + b + 2c) > 7(b + c)(c + a)(a + b) .\]
\end{flushleft}
\subsection*{Page 3}
\begin{flushleft}
\[\int\]
β\par
\bigskip\textbf{Question 3.} (i) Let I n = (sec x + tan x) n dx, where n is a non-negative integer and 0 < β < \pi 2 .\par
\[For n \leq 1, show that\]
\[\sum (\]
1 1\par
\[2 (I n+1 + I n-1) = n (sec β + tan β) n - 1 .\]
Show also that\par
\[\sum (\]
\[I n < n (sec β + tan β) n - 1 .\]
\[\int\]
β\par
\textbf{(ii)} Let J n = (sec x cos β + tan x) n dx, where n is a non-negative integer and 0 < β < \pi 2 .\par
\[For n \leq 1, show that\]
\[\sum (\]
\[J n < n (1 + tan β) n - cos n β .\]
\bigskip\textbf{Question 4.} Let n be a vector of unit length and Π be the plane through the origin perpendicular to n .\par
For any vector x, the projection of x onto the plane Π is defined to be the vector x-(x . n) n .\par
The vectors a and b each have unit length and the angle between them is θ, which satisfies\par
\[0 < θ < \pi . The vector m is given by m = 1 (a + b) .\]
\textbf{(i)} Show that m bisects the angle between a and b.\par
\textbf{(ii)} The vector c also has unit length. The angle between a and c is α, and the angle\par
between b and c is β. Both angles are acute and non-zero.\par
Let a 1 and b 1 be the projections of a and b, respectively, onto the plane through the\par
\[origin perpendicular to c. Show that a 1 . c = 0 and, by considering |a 1 |2 = a 1 . a 1,\]
\[show that |a 1 | = sin α.\]
Show also that the angle φ between a 1 and b 1 satisfies\par
cos θ - cos α cos β\par
\[cos φ = .\]
sin α sin β\par
\textbf{(iii)} Let m 1 be the projection of m onto the plane through the origin perpendicular to c.\par
Show that m 1 bisects the angle between a 1 and b 1 if and only if\par
\[α = β or cos θ = cos(α - β).\]
\end{flushleft}
\subsection*{Page 4}
\begin{flushleft}
\bigskip\textbf{Question 5.} Two curves have polar equations r = a + 2 cos θ and r = 2 + cos 2θ, where r \leq 0 and a is a\par
constant.\par
\textbf{(i)} Show that these curves meet when\par
\[2 cos 2 θ - 2 cos θ + 1 - a = 0 .\]
\[Hence show that these curves touch if a = 1 and find the other two values of a for\]
which the curves touch.\par
\textbf{(ii)} Sketch the curves r = a + 2 cos θ and r = 2 + cos 2θ on the same diagram in the case\par
\[a = 1 . Give the values of r and θ at the points at which the curves touch and justify\]
the other features you show on your sketch.\par
\textbf{(iii)} On two further diagrams, one for each of the other two values of a, sketch both the\par
\[curves r = a + 2 cos θ and r = 2 + cos 2θ. Give the values of r and θ at the points at\]
which the curves touch and justify the other features you show on your sketch.\par
6 (i) For x \ne tan α, the function fα is defined by\par
) [\par
x tan α + 1\par
\[fα(x) = tan -1\]
tan α - x\par
\[where 0 < α < 1 \pi.\]
\[Show that f \leq (x) = .\]
\[α 1 + x^2\]
\[Hence sketch y = fα(x).\]
\[On a separate diagram, sketch y = fα(x) - fβ(x) where 0 < α < β < 1\]
\[\pi.\]
\textbf{(ii)} For 0 \geq x \geq 2\pi and x \ne 1 \pi, 3 \pi, the function g(x) is defined by\par
2 2\par
\[g(x) = tanh -1 (sin x) - sinh -1 (tan x).\]
\[For 1 \pi < x < 3 \pi, show that g \leq (x) = 2 sec x.\]
2 2\par
\[Use this result to sketch y = g(x) for 0 \geq x \geq 2\pi.\]
\end{flushleft}
\subsection*{Page 5}
\begin{flushleft}
\bigskip\textbf{Question 7.} (i) Let\par
iθ iφ\par
e + e\par
\[z = ,\]
eiθ - eiφ\par
\[where θ and φ are real, and θ - φ \ne 2n\pi for any integer n. Show that\]
\[\sum (\]
\[z = i cot 1 (φ - θ)\]
and give expressions for the modulus and argument of z.\par
\textbf{(ii)} The distinct points A and B lie on a circle with radius 1 and centre O. In the complex\par
plane, A and B are represented by the complex numbers a and b, and O is at the\par
\[origin. The point X is represented by the complex number x, where x = a + b and\]
a + b \ne 0. Show that OX is perpendicular to AB.\par
If the distinct points A, B and C in the complex plane, which are represented by the complex\par
numbers a, b and c, lie on a circle with radius 1 and centre O, and h = a + b + c represents\par
the point H, then H is said to be the orthocentre of the triangle ABC.\par
\textbf{(iii)} The distinct points A, B and C lie on a circle with radius 1 and centre O. In the\par
complex plane, A, B and C are represented by the complex numbers a, b and c, and\par
O is at the origin.\par
Show that, if the point H, represented by the complex number h, is the orthocentre\par
\[of the triangle ABC, then either h = a or AH is perpendicular to BC.\]
\textbf{(iv)} The distinct points A, B, C and D (in that order, anticlockwise) all lie on a circle with\par
radius 1 and centre O. The points P , Q, R and S are the orthocentres of the triangles\par
ABC, BCD, CDA and DAB, respectively. By considering the midpoint of AQ, show\par
that there is a single transformation which maps the quadrilateral ABCD on to the\par
quadrilateral QRSP and describe this transformation fully.\par
\end{flushleft}
\subsection*{Page 6}
\begin{flushleft}
\bigskip\textbf{Question 8.} A sequence x 1 , x 2 , . . . of real numbers is defined by x n+1 = x2 n - 2 for n \leq 1 and x 1 = a.\par
\textbf{(i)} Show that if a > 2 then x n \leq 2 + 4 n-1 (a - 2).\par
\textbf{(ii)} Show also that x n \to \infty as n \to \infty if and only if |a| > 2.\par
\textbf{(iii)} When a > 2, a second sequence y 1 , y 2 , . . . is defined by\par
Ax x \cdots x\par
\[y n = 1 x 2 n ,\]
n+1\par
\[where A is a positive constant and n \leq 1.\]
\[Prove that, for a certain value of a, with a > 2, which you should find in terms of A,\]
]\par
\[x^2 - 4\]
n+1\par
\[y n = x\]
n+1\par
\[for all n \leq 1.\]
Determine whether, for this value of a, the second sequence converges.\par
\end{flushleft}
\subsection*{Page 7}
\begin{flushleft}
\bigskip\textbf{Section B: Mechanics}\par
\bigskip\textbf{Question 9.} An equilateral triangle ABC has sides of length a. The points P , Q and R lie on the sides\par
BC, CA and AB, respectively, such that the length BP is x and QR is parallel to CB. Show\par
\[that \sqrt \sqrt\]
\[( 3 cot φ + 1)( 3 cot θ + 1)x = 4(a - x) ,\]
\[where θ = ∠CP Q and φ = ∠BRP .\]
A horizontal triangular frame with sides of length a and vertices A, B and C is fixed on a\par
smooth horizontal table. A small ball is placed at a point P inside the frame, in contact\par
with side BC at a distance x from B. It is struck so that it moves round the triangle P QR\par
described above, bouncing off the frame at Q and then R before returning to point P . The\par
frame is smooth and the coefficient of restitution between the ball and the frame is e.\par
Show that\par
ae\par
\[x = .\]
1 + e\par
Show further that if the ball continues to move round P QR after returning to P , then e = 1.\par
\bigskip\textbf{Question 10.} The origin O of coordinates lies on a smooth horizontal table and the x- and y-axes lie in\par
the plane of the table. A cylinder of radius a is fixed to the table with its axis perpendicular\par
to the x-y plane and passing through O, and with its lower circular end lying on the table.\par
One end, P , of a light inextensible string P Q of length b is attached to the bottom edge of\par
the cylinder at (a, 0). The other end, Q, is attached to a particle of mass m, which rests on\par
the table.\par
Initially P Q is straight and perpendicular to the radius of the cylinder at P , so that Q is at\par
(a, b). The particle is then given a horizontal impulse parallel to the x-axis so that the string\par
immediately begins to wrap around the cylinder. At time t, the part of the string that is\par
\[still straight has rotated through an angle θ, where aθ < b.\]
\textbf{(i)} Obtain the Cartesian coordinates of the particle at this time.\par
Find also an expression for the speed of the particle in terms of θ, θ˙, a and b.\par
\textbf{(ii)} Show that\par
\[θ˙(b - aθ) = u ,\]
where u is the initial speed of the particle.\par
\textbf{(iii)} Show further that the tension in the string at time t is\par
\[mu2\]
\[\sqrt .\]
\[b^2 - 2aut\]
\end{flushleft}
\subsection*{Page 8}
\begin{flushleft}
\bigskip\textbf{Section C: Probability and Statistics}\par
\bigskip\textbf{Question 11.} The continuous random variable X has probability density function\par
\{\par
\[λe -λx for x \leq 0,\]
\[f(x) =\]
0 otherwise,\par
where λ is a positive constant.\par
\[The random variable Y is the greatest integer less than or equal to X, and Z = X - Y .\]
\textbf{(i)} Show that, for any non-negative integer n,\par
\[P(Y = n) = (1 - e -λ )e -nλ.\]
\textbf{(ii)} Show that\par
1 - e -λz\par
\[P(Z < z) = for 0 \geq z \geq 1.\]
1 - e-λ\par
\textbf{(iii)} Evaluate E(Z).\par
\textbf{(iv)} Obtain an expression for\par
\[P(Y = n and z 1 < Z < z 2) ,\]
\[where 0 \geq z 1 < z 2 \geq 1 and n is a non-negative integer.\]
Determine whether Y and Z are independent.\par
\end{flushleft}
\subsection*{Page 9}
\begin{flushleft}
12 (i) In a game, each member of a team of n players rolls a fair six-sided die.\par
The total score of the team is the number of pairs of players rolling the same number.\par
For example, if 7 players roll 3, 3, 3, 3, 6, 6, 2 the total score is 7, as six different pairs\par
of players both score 3 and one pair of players both score 6.\par
Let X ij, for 1 \geq i < j \geq n, be the random variable that takes the value 1 if players i\par
and j roll the same number and the value 0 otherwise.\par
Show that X 12 is independent of X 23 .\par
Hence find the mean and variance of the team's total score.\par
\textbf{(ii)} Show that, if Y i, for 1 \geq i \geq m, are random variables with mean zero, then\par
\}m m\}-1 \}m\par
\[Var(Y 1 + Y 2 + . . . + Y m) = E(Y i 2 ) + 2 E(Y i Y j) .\]
\[i=1 i=1 j=i+1\]
\textbf{(iii)} In a different game, each member of a team of n players rolls a fair six-sided die.\par
The total score of the team is the number of pairs of players rolling the same even\par
number minus the number of pairs of players rolling the same odd number.\par
For example, if 7 players roll 3, 3, 3, 3, 6, 6, 2 the total score is -5.\par
Let Z ij, for 1 \geq i < j \geq n, be the random variable that takes the value 1 if players\par
i and j roll the same even number, the value -1 if players i and j roll the same odd\par
number and the value 0 otherwise.\par
Show that Z 12 is not independent of Z 23 .\par
Find the mean of the team's total score and show that the variance of the team's total\par
\[score is 1 n(n^2 - 1).\]
\end{flushleft}
\end{document}
